%=====================================================================
\section{The RL side by side with the MILP formulation}
\label{sec:compare}

\subsection{Where each MILP variable lives in the RL}

\begin{table}[H]
\centering\small
\caption{Every MILP decision variable and how the same decision appears in an RL episode.}
\begin{tabularx}{\textwidth}{@{}l>{\raggedright\arraybackslash}X>{\raggedright\arraybackslash}p{4.6cm}@{}}
\toprule
MILP & RL equivalent & code \\
\midrule
$y_{ig}=1$ & action $(i,g)$ with $g<G$ was taken & \code{env.res[i] == g} \\
$z_{im}=1$ & action $(i,G+m)$ was taken & \code{env.res[i] == G+m} \\
$x_{uvm}=1$ & cars $u$ and $v$ were added to van $m$ one after the other (or $u=0$ for the first, $v=n+1$ after the last) & route in \code{plan\_of}, sorted by \code{order} \\
$u_m=1$ & van $m$ got at least one car & \code{van\_cnt[m] > 0} \\
$o_{ij}=1$ & on a shared gun, car $i$ was added before car $j$ & \code{order[i] < order[j]} \\
$S_i$ & earliest start, computed when $i$ is added & \code{env.S[i]} (from \code{start\_g} / \code{start\_v}) \\
$q^\Fcs_i$, $q^\Mcs_i$ & queue / van wait at that moment & \code{wait\_g}, \code{wait\_v} \\
$C_i$ & $S_i$ plus the charging time & \code{env.C[i]} \\
\bottomrule
\end{tabularx}
\end{table}

The RL never creates these variables; they are \emph{recovered} from the order and the choices of the episode. The MILP has one binary per possible pair of cars per van and per shared gun (thousands); the RL instead makes $n$ decisions, each among at most $n\times K$ pairs.

\subsection{Where each MILP constraint lives in the RL}

{\small
\begin{longtable}{@{}l>{\raggedright\arraybackslash}p{4.2cm}>{\raggedright\arraybackslash}p{10.0cm}@{}}
\toprule
no. & MILP constraint & how the RL guarantees it \\
\midrule
\endhead
(2) & one resource per car & each car is chosen once and then marked \code{done}; masks exclude done cars (\code{options}, line 192) \\
(3), (20) & start after arrival / request & by construction: $S_i=\max(r_i+t_{if},a_g)$ or $\max(r_i,a_m+t_{\ell_mi})$ \\
(4) & definition of the waits & computed directly: \code{wait\_g}, \code{wait\_v} \\
(5) & queue $\le W_i$ & gun mask: \code{wait\_g <= W} (line 170) \\
(6) & only van users wait for vans & automatic: a station user has no van wait \\
(7)--(8) & no overlap on a gun & by construction: a car is added at the end of the gun's queue and starts after $a_g$, the completion of the previous car \\
(9) & gun-load cut & not needed: it is a \emph{valid inequality} (implied by the others); the MILP has it only to help Gurobi's bounds \\
(10) & gun symmetry & symmetry mask: of identical idle guns of one station only the first is offered (lines 169, 171) \\
(11)--(14) & route structure (leave, visit once, return) & by construction: a van's cars form one sequence from the depot; the return is implied \\
(15) & at most $Q$ cars & van mask: $\kappa_m<Q$ (line 185) \\
(16) & battery & van mask: $E_i/\eta+\beta(d_{\ell_mi}+d_{i0})\le b_m$ (lines 182, 186); $b_m$ decreases after each car \\
(17)--(18) & van timing, one car at a time & by construction: $S_i\ge a_m+t_{\ell_mi}$, where $a_m$ is the completion of the van's previous car \\
(19) & back at the depot by $T$ & van mask: $C_i+t_{i0}\le T$ (line 187); checked for the newest last car each time \\
(21) & charged by $T$ & masks: $C_i\le T$ for guns (line 170); for vans implied by (19) \\
(22) & van symmetry & symmetry mask: only the first unused van is offered (lines 183--185) \\
(23) & domains & automatic \\
\bottomrule
\end{longtable}}

\textbf{One difference in what ``infeasible'' means.} The MILP either finds a plan satisfying all rules or reports none. The RL always produces a complete assignment, but if a car was dropped, \code{plan\_of} puts it on the resource that finishes it earliest and breaks a rule; the checker marks the plan invalid. So ``no valid plan'' is reported the same way for both, and an invalid RL plan is never counted as a solution.

\subsection{The objective}

Proposition~\ref{prop:sum} shows that the sum of the RL's step costs is exactly $Z$ of (1). The split of the six terms across the steps is in the boxed formulas of Section~\ref{sec:reward}. In practice the RL's $Z$ and the checker's $Z$ are compared for every solved instance (\code{solve\_rl\_f4.py}, line 130), and a warning is written if they differ by more than $10^{-6}$ relative; in the server run this warning appears in none of the 31 RL runs ($31\times560$ instances).

\subsection{What each method can and cannot do}

\begin{table}[H]
\centering\small
\begin{tabularx}{\textwidth}{@{}>{\raggedright\arraybackslash}p{3.4cm}>{\raggedright\arraybackslash}X>{\raggedright\arraybackslash}X@{}}
\toprule
 & MILP & RL \\
\midrule
how decisions are made & all together, by branch and bound over the binary variables & one at a time; earlier decisions are never changed \\
timing & optimised together with everything else & earliest start (optimal for the chosen orders, Proposition~\ref{prop:timing}) \\
rules & constraints (2)--(23), with big-$M$ terms & masks + construction (table above) \\
quality guarantee & proven optimal, or a known gap & none; judged against the MILP \\
time per instance & seconds to 300 s & about 0.1 s (520 tries) \\
when it fails & no plan found in time & a car left without allowed resource \\
learning & none; each instance from scratch & trained once on random instances, reused \\
\bottomrule
\end{tabularx}
\end{table}

\inshort{Every MILP rule is either a mask (range, queue limit, seats, battery, deadline, back by $T$), a consequence of adding cars at the end of sequences with earliest start times (no overlap, route structure, van timing), or not needed (the gun-load cut). The rewards add up to exactly the same $Z$.}
